% !TEX root = ../main.tex
% Ouverture : titre et ce que fait le moteur

% ======================================================================= titre
\begin{frame}[plain]
  \titlepage
\end{frame}

% ============================================================= ce que fait le moteur
\begin{frame}{Ce que fait le moteur}
  \framesubtitle{des particules, des ressorts, des forces}
  \begin{columns}[c,onlytextwidth]
    \begin{column}{0.58\textwidth}
      \begin{diagramoverlay}[width=\linewidth]{figures/scenes/scene_05.png}
        \astpoint{0.737}{0.94}{0.79}{0.80}{west}{ancre\\épinglée}
        \astpoint{0.722}{0.78}{0.79}{0.60}{west}{ressort\\étiré}
        \astpoint{0.5}{0.32}{0.5}{0.10}{north}{400 particules\astsep 1\,482 ressorts}
      \end{diagramoverlay}
    \end{column}
    \begin{column}{0.38\textwidth}
      \astbig{13}{scènes, du pendule au tissu}
      \vspace{1mm}
      \astbig[phelmared]{4}{forces : poids, frottement, ressort, vent}
      \vspace{1mm}
      \astbig[timagreen]{1/120\,s}{pas fixe, Euler ou Verlet en direct}
    \end{column}
  \end{columns}
  \astcue{
    \cue{montrer le tissu en direct, \cuekey{attraper une particule}}
    \cuenext{comment c'est organisé}
  }
\end{frame}
